\subsection{Coefficients of a Module}

Let E be a left A-module. Let \(E^{*}\) be the dual of the vector space E over K; we endow it with its natural right A-module structure. Given elements x of E and \(x^{*}\) of \(E^{*}\) , we denote by \(c_{\mathrm{E}}(x, x^{*})\) the linear form

\[
c _ {\mathrm{E}} (x, x ^ {*}): a \mapsto \langle x ^ {*}, a x \rangle\tag{5}
\]

on A. These linear forms are called the coefficients of E. The subspace of \(\mathrm{A}^*\) generated by the coefficients of E is denoted by \(\Theta_{\mathrm{E}}(\mathrm{A})\) . It is nonzero if E is nonzero.

If F is an A-module isomorphic to E, then E and F have the same coefficients and we have \(\Theta_{\mathrm{F}}(\mathrm{A}) = \Theta_{\mathrm{E}}(\mathrm{A})\) . Let us view \(E^{*}\) as a left \(A^{\circ}\) -module; a coefficient of E is a coefficient of \(E^{*}\) , and \(\Theta_{\mathrm{E}}(\mathrm{A}) \subset \Theta_{\mathrm{E}^{*}}(\mathrm{A}^{\circ})\) . Consequently, if E is finite-dimensional over K, then E and \(E^{*}\) have the same coefficients, and we have \(\Theta_{\mathrm{E}}(\mathrm{A}) = \Theta_{\mathrm{E}^{*}}(\mathrm{A}^{\circ})\) .

Let \((\mathrm{M},\pi)\) be a linear representation of the algebra A. The coefficients of the A-module \(\mathrm{M}_{\pi}\) are also called the coefficients of the representation \(\pi\) . The coefficient \(c_{\mathrm{M}_{\pi}}(x,x^{*})\) , for \(x\in \mathrm{M}\) and \(x^{*}\in \mathrm{M}^{*}\) , is also denoted by \(c_{\pi}(x,x^{*})\) ; the vector space \(\Theta_{\mathrm{M}_{\pi}}(\mathrm{A})\) is also denoted by \(\Theta_{\pi}(\mathrm{A})\) .

Remarks. --- 1) Suppose that M is finite-dimensional over K. Let \((e_1, \ldots, e_n)\) be a basis of V, and let \((e_1^*, \ldots, e_n^*)\) be its dual basis. Denote by \((\pi_{ij}(a))\) the matrix of \(\pi(a)\) with respect to the basis \((e_1, \ldots, e_n)\) of V; we have \(\pi_{ij} = c_\pi(e_j, e_i^*)\) . The mapping \(a \mapsto (\pi_{ij}(a))\) is a homomorphism from A to \(\mathbf{M}_n(\mathrm{K})\) ; such a homomorphism is sometimes called a matrix representation of the algebra A.

\begin{enumerate}
\def\labelenumi{\arabic{enumi})}
\setcounter{enumi}{1}
\tightlist
\item
  Let E be an A-module, and let \(E'\) be a submodule of E. By the corollary of Theorem 5 of II, §7, No.~5, p.~299, we have \(\Theta_{\mathrm{E}'}(\mathrm{A}) \subset \Theta_{\mathrm{E}}(\mathrm{A})\) and \(\Theta_{\mathrm{E}/\mathrm{E}'}(\mathrm{A}) \subset \Theta_{\mathrm{E}}(\mathrm{A})\) .
\end{enumerate}

Let \(\gamma_{\mathrm{E}}\) be the unique K-linear mapping from \(\mathrm{E} \otimes_{\mathrm{K}} \mathrm{E}^{*}\) to \(\mathrm{A}^{*}\) that sends \(x \otimes x^{*}\) to \(c_{\mathrm{E}}(x, x^{*})\) . It is (A,A)-bilinear, and its image is \(\Theta_{\mathrm{E}}(\mathrm{A})\) . We deduce from it A-linear mappings \(c_{\mathrm{E}}^{\prime}: \mathrm{E} \to \operatorname{Hom}_{\mathrm{A}}(\mathrm{E}^{*}, \mathrm{A}^{*})\) and \(c_{\mathrm{E}}^{\prime \prime}: \mathrm{E}^{*} \to \operatorname{Hom}_{\mathrm{A}}(\mathrm{E}, \mathrm{A}^{*})\) such that

\[
c _ {\mathrm{E}} ^ {\prime} (x) (x ^ {*}) = c _ {\mathrm{E}} (x, x ^ {*}) = c _ {\mathrm{E}} ^ {\prime \prime} (x ^ {*}) (x).
\]

Denote by \(\theta_{E}\) the K-linear mapping \(\theta_{E}: E \otimes E^{*} \to \operatorname{End}_{K}(E)\) characterized by the relation

\[
\theta_ {\mathrm{E}} (x \otimes x ^ {*}) (y) = \langle x ^ {*}, y \rangle x
\]

for \(x, y \in E\) and \(x^{*} \in E^{*}\) . Its image is the set \(\operatorname{End}_{\mathrm{K}}^{\mathrm{f}}(\mathrm{E})\) of endomorphisms of E of finite rank (VIII, p.~463). By the definition of the trace (loc. cit.), we have

\[
\operatorname{Tr} \left(\theta_ {\mathrm{E}} \left(x \otimes x ^ {*}\right)\right) = \langle x ^ {*}, x \rangle
\]

for \(x \in E\) and \(x^{*} \in E^{*}\) ; we therefore have

\[
\langle \gamma_ {\mathrm{E}} (x \otimes x ^ {*}), a \rangle = \langle x ^ {*}, a x \rangle = \operatorname{Tr} (\theta_ {\mathrm{E}} (a x \otimes x ^ {*})) = \operatorname{Tr} (\theta_ {\mathrm{E}} (x \otimes x ^ {*}) a).
\]

This gives the relation

\[
\langle \gamma_ {\mathrm{E}} (h), a \rangle = \operatorname{Tr} (\theta_ {\mathrm{E}} (h) \circ a _ {\mathrm{E}})
\]

for \(a \in \mathrm{A}\) and \(h \in \mathrm{E} \otimes_{\mathrm{K}} \mathrm{E}^*\) . This proves that \(\Theta_{\mathrm{E}}(\mathrm{A})\) is the set of linear forms \(a \mapsto \operatorname{Tr}(u \circ a_{\mathrm{E}})\) on \(\mathrm{A}\) , where \(u\) runs through \(\operatorname{End}_{\mathrm{K}}^{\mathrm{f}}(\mathrm{E})\) .

Lemma 1. --- The mapping \(c_{\mathrm{E}}^{\prime \prime}\) is bijective. If \(\mathrm{E}\) is finite-dimensional, then so is \(c_{\mathrm{E}}^{\prime}\) .

When F is a right A-module and G is a K-vector space, we defined in II, §4, No.~1, p.~268, Proposition 1, a) an isomorphism \(\gamma\) of K-vector spaces from \(\mathrm{Hom}_{\mathrm{K}}(\mathrm{F} \otimes_{\mathrm{A}} \mathrm{E}, \mathrm{G})\) to \(\mathrm{Hom}_{\mathrm{A}}(\mathrm{E}, \mathrm{Hom}_{\mathrm{K}}(\mathrm{F}, \mathrm{G}))\) . This isomorphism sends \(\varphi : \mathrm{F} \otimes_{\mathrm{A}} \mathrm{E} \to \mathrm{G}\) to the homomorphism \(e \mapsto (f \mapsto \varphi(f \otimes e))\) . Through the canonical isomorphism from \(\mathrm{A}_d \otimes_{\mathrm{A}} \mathrm{E}\) to E (II, §3, No.~4, p.~249), \(\gamma\) can be identified with \(c_{\mathrm{E}}''\) when \(\mathrm{F} = \mathrm{A}_d\) and \(\mathrm{G} = \mathrm{K}\) .

Analogously, the isomorphism \(\beta\) defined in II, §4, No.~1, p.~268, Proposition 1, b) specializes to an isomorphism \(\beta\) from \(\mathrm{Hom}_{\mathrm{K}}(\mathrm{E}^{*}\otimes_{\mathrm{A}}\mathrm{A}_{s},\mathrm{K})\) to \(\mathrm{Hom}_{\mathrm{A}}(\mathrm{E}^{*},\mathrm{Hom}_{\mathrm{K}}(\mathrm{A}_{s},\mathrm{K}))\) . When E is finite-dimensional, E can be canonically identified with \(\mathrm{Hom}_{\mathrm{K}}(\mathrm{E}^{*},\mathrm{K})\) and \(\mathrm{E}^{*}\) with \(\mathrm{E}^{*}\otimes_{\mathrm{A}}\mathrm{A}_{s}\) ; the homomorphism \(\beta\) is then identified with \(c_{\mathrm{E}}^{\prime}\) .

PROPOSITION 2. --- Let E be a left A-module.

\begin{enumerate}
\def\labelenumi{\alph{enumi})}
\item
  The set of coefficients of E is the union of the images of the A-linear mappings from E to A*.
\item
  Suppose, moreover, that E has finite dimension n over K. Then the left A-module \(\Theta_{\mathrm{E}}(\mathrm{A})\) is isomorphic to a quotient of \(E^{n}\) , the A-module E is isomorphic to a submodule of \(\Theta_{\mathrm{E}}(\mathrm{A})^{n}\) , and every element of \(\Theta_{\mathrm{E}}(\mathrm{A})\) is a coefficient of \(E^{n}\) .
\end{enumerate}

Assertion a) follows from the surjectivity of \(c_{\mathrm{E}}^{\prime \prime}\) .

Let us prove b). Let \((e_1^*,\ldots ,e_n^*)\) be a basis of \(\mathrm{E}^*\) over K. Since \(\Theta_{\mathrm{E}}(\mathrm{A})\) is generated by the coefficients of E, the A-linear mapping

\[
(x _ {1}, \ldots , x _ {n}) \mapsto \sum_ {i = 1} ^ {n} c _ {\mathrm{E}} (x _ {i}, e _ {i} ^ {*})
\]

from \(\mathrm{E}^n\) to \(\Theta_{\mathrm{E}}(\mathrm{A})\) is surjective. By a), every element of \(\Theta_{\mathrm{E}}(\mathrm{A})\) is a coefficient of \(\mathrm{E}^n\) . Moreover, the A-linear mapping

\[
x \mapsto \left(c _ {\mathrm{E}} (x, e _ {1} ^ {*}), \ldots , c _ {\mathrm{E}} (x, e _ {n} ^ {*})\right)
\]

from E to \(\Theta_{\mathrm{E}}(\mathrm{A})^{n}\) is injective; this gives b).

Remark 3. --- Let E be a left A-submodule of \(\mathrm{A}^*\) . Let \(\varepsilon\) be the linear form \(y \mapsto y(1)\) on E. For every \(x\) in E, we have \(x = c_{\mathrm{E}}(x, \varepsilon)\) , so E is an A-submodule of \(\Theta_{\mathrm{E}}(\mathrm{A})\) .

